Define \(s_{0,B}^2\) and \(\mathsf A_{m,B}\) as in the statement.  The pilot-oracle-regret theorem and \(r_{V,m}=o(1)\) give
\[
 \max_{d\in\mathcal D}
 |\widehat\alpha_{d,m}-\alpha_d^*(P)|=o_P(1).               \tag{1}
\]
Because \(\mathcal D\) is finite and all costs and source variances are positive and finite, there is \(a_0\in(0,1/2)\) such that
\[
 2a_0\le \alpha_d^*(P)\le1-2a_0,
 \qquad d\in\mathcal D.                                    \tag{2}
\]
Equations (1)--(2) imply that, with probability tending to one, the selected share lies in \([a_0,1-a_0]\).  Hence both selected counts are positive and proportional to \(B\), so
\[
 P(\mathsf A_{m,B})\longrightarrow1.                        \tag{3}
\]
The uniform floor-count conclusion of the pilot-oracle-regret theorem, evaluated at the random member \(\widehat d_m\) of the finite family, yields
\[
 s_{0,B}^2-\mathcal V_{\widehat d_m}^*(P)=o_P(1).            \tag{4}
\]
Every expression below is evaluated on \(\mathsf A_{m,B}\); the fixed measurable extensions on its complement do not affect any limit by (3).

By the definition of \(\widehat s_B^2\),
\[
\begin{aligned}
 |\widehat s_B^2-s_{0,B}^2|
 &\le \frac{B}{N_{E,B}}
 |\widetilde V_{E,\widehat d_m,B}-V_{E,\widehat d_m}|\\
 &\quad+\frac{B}{N_{O,B}}
 |\widetilde V_{O,\widehat d_m,B}-V_{O,\widehat d_m}|.
\end{aligned}                                                \tag{5}
\]
The count ratios in (5) are \(O_P(1)\) by (1)--(2), while the uniform main-study variance error is \(O_P(r_{\widetilde V,B})=o_P(1)\).  Therefore
\[
 \widehat s_B^2-s_{0,B}^2=o_P(1),
 \qquad
 \widehat s_B^2-\mathcal V_{\widehat d_m}^*(P)=o_P(1).      \tag{6}
\]
Nondegeneracy and finiteness further give
\[
 0<\min_{d\in\mathcal D}\mathcal V_d^*(P)
 \le \max_{d\in\mathcal D}\mathcal V_d^*(P)<\infty.      \tag{7}
\]
It follows from (4), (6), and (7) that \(\widehat s_B/s_{0,B}\to_P1\).

We next prove the oracle-studentized limit conditional on any deterministic pilot path.  Fix deterministic sequences \(d_B\in\mathcal D\) and \(\alpha_B\in(0,1)\) such that
\[
 |\alpha_B-\alpha_{d_B}^*(P)|\longrightarrow0.              \tag{8}
\]
Let \(e_B,o_B,n_B\) be the associated counts in the statement.  From any subsequence, finiteness of \(\mathcal D\) permits extraction of a further subsequence on which \(d_B=d\) is constant.  Along this further subsequence,
\[
 \frac{e_B}{B}\longrightarrow\frac{\alpha_d^*(P)}{c_{E,d}},
 \qquad
 \frac{o_B}{B}\longrightarrow\frac{1-\alpha_d^*(P)}{c_{O,d}}. \tag{9}
\]
Thus \(e_B,o_B\to\infty\), \(n_B\asymp B\), and
\[
 \lambda_d:=\lim_{B\to\infty}\frac{e_B}{o_B}
 =\frac{\alpha_d^*(P)c_{O,d}}
 {(1-\alpha_d^*(P))c_{E,d}}\in(0,\infty).                  \tag{10}
\]

We now map the fixed-design scores exactly.  For \(g\in\{E,O\}\), write \(p_{g,1}=P(A=1\mid G=g)\) and \(p_{g,0}=1-p_{g,1}\).  The arm-specific treatment normalization in the difference of the two conditional-mean estimators is the identity
\[
 \frac{A-p_{g,1}}{p_{g,1}p_{g,0}}R_g(A)
 =\frac{\mathbf 1\{A=1\}}{p_{g,1}}R_g(1)
  -\frac{\mathbf 1\{A=0\}}{p_{g,0}}R_g(0).                 \tag{11}
\]
For the experimental source take
\[
 R_E(a)=h_d(S_{3d},S_{2d},a,X)-\mu_{d,a},
\]
and for the observational source take
\[
 R_O(a)=q_d(S_{2d},S_{1d},a,X)
 \{Y-h_d(S_{3d},S_{2d},a,X)\}.
\]
Equation (11) is precisely \(\phi_{E,d}\) and \(\phi_{O,d}\), respectively.  Consequently the source-score second moments in the cited IKM fixed-design theorem are exactly \(V_{E,d}\) and \(V_{O,d}\); no change of treatment normalization is hidden in the notation.

The remaining premises of the cited theorem also map directly.  The fold count \(L\) is fixed.  Conditional on the pilot, the two source samples are mutually independent and independent and identically distributed within source.  The theorem's causal-subclass premise supplies both identification routes, both latent bridges, and the observed bridge equations, with the nuisance limits equal to \(h_d\) and \(q_d\).  The four sequence-wise rate assumptions are exactly IKM Assumption~8 in \(L_2(P_{B,d,\alpha_B}^{\circ})\), the norm of its artificial two-source superpopulation.  In particular, no observational-source norm is substituted for that norm.  From \(n_B\asymp B\),
\[
 \min\{\delta_{h,B}\rho_{q,B},\rho_{h,B}\delta_{q,B}\}
 =o(B^{-1/2})=o(n_B^{-1/2}).                                \tag{12}
\]
Finally, the nondegeneracy assumption supplies finite, strictly positive source-score variances.  The cited fixed-design theorem therefore gives
\[
 \sqrt{n_B}(\widehat\tau_B-\tau)
 \rightsquigarrow
 \mathcal N\!\left(0,
 \frac{1+\lambda_d}{\lambda_d}V_{E,d}
 +(1+\lambda_d)V_{O,d}\right).                             \tag{13}
\]
The variance in (13) is exactly the count-normalized two-source variance, because
\[
\begin{aligned}
 n_B\left(\frac{V_{E,d}}{e_B}+\frac{V_{O,d}}{o_B}\right)
 &=\left(1+\frac{o_B}{e_B}\right)V_{E,d}
   +\left(1+\frac{e_B}{o_B}\right)V_{O,d}\\
 &\longrightarrow
 \frac{1+\lambda_d}{\lambda_d}V_{E,d}
 +(1+\lambda_d)V_{O,d}.                                    
\end{aligned}                                                \tag{14}
\]
Combining (13)--(14) gives, along every deterministic sequence satisfying (8),
\[
 \frac{\sqrt B(\widehat\tau_B-\tau)}
 {\{B(V_{E,d}/e_B+V_{O,d}/o_B)\}^{1/2}}
 \rightsquigarrow\mathcal N(0,1).                          \tag{15}
\]

It remains to pass from deterministic pilot paths to the random pilot without assuming a uniform limit theorem.  From (1), choose a deterministic sequence \(a_m\downarrow0\) such that
\[
 P\left\{
 \max_{d\in\mathcal D}
 |\widehat\alpha_{d,m}-\alpha_d^*(P)|>a_m\right\}
 \longrightarrow0,                                         \tag{16}
\]
and denote the complementary event by \(\mathsf H_m\).  Let \(d_{\mathrm{BL}}\) be bounded-Lipschitz distance.  We claim that the conditional law, given any pilot realization in \(\mathsf H_m\), of the oracle-studentized statistic in (15) converges uniformly over those realizations to \(\mathcal N(0,1)\).  If not, there would be a positive constant, a subsequence, and a violating deterministic pilot realization at each index.  Their selected designs and shares satisfy (8); extracting a constant-design subsequence and applying (9)--(15) forces the bounded-Lipschitz distance to zero, a contradiction.  Thus this uniformity is derived from the every-deterministic-sequence premise rather than assumed.

Pilot independence makes the conditional main-study distribution given each pilot realization exactly the fixed-design distribution just analyzed.  Integrating the conditional bounded-Lipschitz distances over \(\mathsf H_m\), and bounding the contribution of \(\mathsf H_m^c\) by (16), proves
\[
 \frac{\sqrt B(\widehat\tau_B-\tau)}{s_{0,B}}
 \rightsquigarrow\mathcal N(0,1)                           \tag{17}
\]
under the joint pilot and main-study law.  Equations (6)--(7), (17), and Slutsky's theorem yield
\[
 \frac{\sqrt B(\widehat\tau_B-\tau)}{\widehat s_B}
 \rightsquigarrow\mathcal N(0,1).                          \tag{18}
\]
Finally, on \(\mathsf A_{m,B}\),
\[
 \{\tau\in I_B\}
 =\left\{
 \left|\frac{\sqrt B(\widehat\tau_B-\tau)}{\widehat s_B}\right|
 \le z_{1-\eta/2}\right\}.                                \tag{19}
\]
The standard normal law is continuous at \(\pm z_{1-\eta/2}\).  Equations (3), (18), and (19) therefore imply \(P\{\tau\in I_B\}\to1-\eta\), as claimed.